% !PRINTED 449
\weberpage{449}
Between these three poles, however, there is an essential difference. For if we apply the substitution $\omega$ to them, the first of them remains unchanged. It is therefore at least six-fold and belongs to the substitutions
\[
1, \chi, \chi^2, \omega, \omega\chi, \omega\chi^2.
\]
The two other poles, namely
\weq{2}{x_1 = \varrho x_2 = \varrho^2 x_3, \quad x_1 = \varrho^2 x_2 = \varrho x_3,}
pass into one another under the substitution $\omega$, and we must show that they are altered by every other substitution of the group except $\chi, \chi^2$, and hence that they are only three-fold. If for the moment we denote the two points (2) by $\pi$ and $\pi'$, we can first say that all substitutions of $G_{168}$ which leave $\pi$ unchanged must form a group $G_\pi$,
and indeed a divisor of $G_{168}$; the group $G_\pi$ in turn contains the cyclic group 1, $\chi$, $\chi^2$.

If $\sigma$ is any substitution of $G_\pi$, then $\pi$ occurs among the poles of $\sigma$. From this it follows that $\sigma$ can be neither of degree $7$ nor of degree $4$. For the coordinates of the poles of these two degrees are rational functions of $\varepsilon, i$, while the coordinates of $\pi$ contain the irrationality $\varrho$, which is independent of these (Bd. I, \S.\,134, cf. the supplement at the end of this volume). But that $\sigma$ also cannot be of degree $2$ follows from the fact that $\pi$ lies on none of the principal axes. For if it lay on a principal axis, then, because $\varrho$ cannot be expressed rationally in terms of $\varepsilon$, as is to be derived from (2) by means of the equation
\[
1 + \varrho + \varrho^2 = 0
\]
one of the 21 relations of \S.\,118, (7) would have to hold, namely
\[
\begin{aligned}
\alpha &= \beta\varepsilon^{r+3} &&= \gamma\varepsilon^{-2r+1}, \\
\alpha &= \beta\varepsilon^{-3r+3} &&= \gamma\varepsilon^{-r+1}, \\
\alpha &= \beta\varepsilon^{2r+3} &&= \gamma\varepsilon^{3r+1},
\end{aligned}
\]
of which evidently none is possible. The group $G_\pi$ therefore contains, besides the identical substitution, only substitutions of degree $3$, and the degree of $G_\pi$ must be a power of 3. But since 168 is not divisible by 9, the degree of $G_\pi$ must be equal to 3. From this it follows:

5. \gesperrt{There are 56 three-fold poles, which belong two by two to the same substitution, and accordingly arrange themselves in 28 pairs. Each of these poles}
%CONT>>
